%%═══════════════════════════════════════════════════════════════════
%% Lecture 11 — Nuclear Energy: Fusion, Fission & Reactors
%% 77603 — Nuclear Physics
%% Date: 23 June 2026
%%═══════════════════════════════════════════════════════════════════

\section{Lecture 11 — Nuclear Energy: Fusion, Fission \& Reactors}
\label{sec:lec11}

\noindent\textbf{\textcolor{doc}{Lecture date: 23 June 2026}}

%%───────────────────────────────────────────────────────────────────
\subsection*{Overview}
%%───────────────────────────────────────────────────────────────────
This lecture asks how nuclear binding energy becomes usable macroscopic energy.
The unifying picture is the binding-energy curve: nuclei below the iron/nickel
region can usually release energy by \emph{fusion}, while very heavy nuclei can
release energy by \emph{fission}. The two processes look opposite---joining
light nuclei versus splitting heavy nuclei---but the accounting is the same: if
the final nuclear configuration has larger total binding energy than the initial
configuration, the missing rest mass appears as kinetic energy, photons,
neutrinos, and ultimately heat.

The lecture then develops three linked topics.
\begin{itemize}[nosep]
  \item \textbf{Fusion}: Coulomb-barrier penetration, the Gamow window in
        stars, magnetic/inertial confinement on Earth, and the distinction
        between physical and engineering break-even.
  \item \textbf{Fission}: the liquid-drop competition between surface tension
        and Coulomb repulsion, induced fission of uranium, prompt neutrons,
        fission-fragment distributions, and why \nucleus{U}{235} and
        \nucleus{U}{238} behave so differently.
  \item \textbf{Reactors}: neutron multiplication, the four-factor and
        six-factor formulas, criticality, delayed neutrons, moderators,
        coolants, control rods, enrichment, natural reactors, and the main
        engineering/environmental issues of nuclear power.
\end{itemize}

%%═══════════════════════════════════════════════════════════════════
\subsection{The Binding-Energy Curve as an Energy Map}
\label{subsec:lec11:binding_curve_energy}
%%═══════════════════════════════════════════════════════════════════

The starting point is the familiar plot of binding energy per nucleon,
$B/A$, versus mass number $A$. It rises rapidly from hydrogen to light nuclei,
continues upward more slowly, reaches a broad maximum around
$A\simeq 56$--$60$ (iron/nickel region), and then decreases slowly for very
heavy nuclei.

\dfn[dfn:lec11:energy_release]{Nuclear Energy Release}{%
For a reaction
\begin{equation}
  \sum_i X_i \longrightarrow \sum_f Y_f,
\end{equation}
the energy release is
\begin{equation}
  \boxed{Q=\left(\sum_i M_i-\sum_f M_f\right)c^2
        = B_f^{\rm tot}-B_i^{\rm tot}.}
\end{equation}
The reaction releases energy when the total final binding energy is larger than
the total initial binding energy.}

\paragraph{Physical intuition.}
Binding energy is negative potential energy written as a positive number: a more
strongly bound final state has lower rest mass. Thus energy is released whenever
nucleons rearrange themselves into a more tightly bound configuration. For light
nuclei, moving rightward/upward on the $B/A$ curve usually means fusion releases
energy. For very heavy nuclei, moving leftward toward two medium-mass fragments
also raises the total binding energy per nucleon, so fission releases energy.

\nt{The important comparison is not the binding energy of one nucleus alone, but
\emph{the total binding energy of all nucleons before and after}. Even if $Z$
and $N$ are separately conserved, rearranging nucleons among different nuclei can
change the total binding energy by many MeV per nucleon.}

%%═══════════════════════════════════════════════════════════════════
\subsection{Fusion: Why It Releases Energy and Why It Is Hard}
\label{subsec:lec11:fusion_basics}
%%═══════════════════════════════════════════════════════════════════

\dfn[dfn:lec11:fusion]{Fusion}{%
Fusion is a reaction in which two light nuclei combine into a heavier nuclear
system,
\begin{equation}
  X_1+X_2\longrightarrow Y+\text{energy/products}.
\end{equation}
For $A\lesssim 60$, fusion can increase the total binding energy and therefore
release energy.}

A schematic first step in the proton--proton chain is
\begin{equation}
  p+p\longrightarrow d+e^+ + \nu_e,
  \label{eq:lec11:pp_first_step}
\end{equation}
where one proton converts to a neutron through the weak interaction. Charge is
conserved because the positron carries positive charge; lepton number is
conserved by the emitted neutrino. The positron later slows down in the plasma
and annihilates with an electron,
\begin{equation}
  e^+ + e^- \longrightarrow 2\gamma.
\end{equation}
The neutrino usually escapes, carrying away energy that is essentially lost from
the point of view of heat production.

\subsubsection*{The Coulomb Barrier}
The reason fusion is difficult is not the nuclear force. Once two nuclei are
within a few femtometres, the strong interaction can bind them. The difficulty is
bringing two positively charged nuclei close enough to reach that range. Outside
the nuclear-force region the potential is approximately
\begin{equation}
  V_C(r)=\frac{1}{4\pi\epsilon_0}\frac{Z_1Z_2e^2}{r}.
\end{equation}
The larger $Z_1Z_2$ is, the higher and thicker the Coulomb barrier becomes.
This is why proton--proton or deuterium--tritium fusion is much easier than, say,
carbon--carbon fusion.

Classically, a particle with energy below the barrier turns around. Quantum
mechanically, it can tunnel. In WKB form the tunnelling probability is
\begin{equation}
  P(E)\sim e^{-2G(E)},
  \qquad
  G(E)=\frac{1}{\hbar}\int_{r_1}^{r_2}\sqrt{2\mu\,[V_C(r)-E]}\,\dd r,
  \label{eq:lec11:fusion_wkb}
\end{equation}
where $r_1,r_2$ are the classical turning points and $\mu$ is the reduced mass.
For a Coulomb barrier this leads to the standard Gamow behaviour
\begin{equation}
  \boxed{P(E)\propto \exp\left[-\sqrt{\frac{E_G}{E}}\right]},
  \label{eq:lec11:gamow_probability}
\end{equation}
with $E_G\propto \mu (Z_1Z_2e^2)^2/\hbar^2$. The exact prefactor is less
important here than the lesson: the probability rises extremely rapidly with
energy, because the forbidden barrier width shrinks.

%%═══════════════════════════════════════════════════════════════════
\subsection{The Gamow Window in Stars}
\label{subsec:lec11:gamow_window}
%%═══════════════════════════════════════════════════════════════════

In a stellar plasma the particles do not all have the same energy. Their kinetic
energies are approximately Maxwell--Boltzmann distributed,
\begin{equation}
  \frac{\dd n}{\dd E}\propto \sqrt{E}\,e^{-E/k_BT}.
\end{equation}
Low-energy particles are abundant but tunnel poorly; high-energy particles
tunnel well but are exponentially rare. The reaction rate is controlled by the
product of these two factors:
\begin{equation}
  \text{rate integrand}\propto e^{-E/k_BT}\,
  e^{-\sqrt{E_G/E}}
  \quad\text{(up to slower factors).}
  \label{eq:lec11:gamow_integrand}
\end{equation}

\thm[th:lec11:gamow_window]{Gamow Window}{%
The dominant stellar fusion reactions occur in a narrow energy interval above
$k_BT$, where the falling Maxwell tail and the rising tunnelling probability
balance. This interval is called the \emph{Gamow window}.}

\subsubsection*{Derivation of the Peak Energy}
To estimate the peak, minimize the exponent
\begin{equation}
  \Phi(E)=\frac{E}{k_BT}+\sqrt{\frac{E_G}{E}}.
\end{equation}
Then
\begin{equation}
  \dv{\Phi}{E}
  =\frac{1}{k_BT}-\frac{1}{2}\sqrt{E_G}\,E^{-3/2}=0.
\end{equation}
Therefore
\begin{equation}
  E_0^{3/2}=\frac{1}{2}\sqrt{E_G}\,k_BT,
\end{equation}
and hence
\begin{equation}
  \boxed{E_0=\left[\frac{(k_BT)^2E_G}{4}\right]^{1/3}.}
  \label{eq:lec11:gamow_peak}
\end{equation}
This $E_0$ is typically several times larger than $k_BT$. Thus the relevant
fusion reactions in the Sun do not occur at the thermal average energy itself,
but in the high-energy tail of the plasma distribution.

\paragraph{Stellar intuition.}
Stars succeed where laboratory devices struggle because they have enormous mass,
gravity, pressure, and time. The Sun's core temperature is of order
$10^7\,\mathrm K$, lower than many terrestrial fusion designs, but solar material
is gravitationally confined for billions of years. Each particle gets many
chances to tunnel.

\paragraph{Beyond hydrogen burning.}
As stars evolve, they may burn helium, carbon, oxygen, neon, and silicon,
eventually producing nuclei near iron/nickel. Each later burning stage requires
higher temperature because larger $Z_1Z_2$ raises the Coulomb barrier. Once the
core reaches the iron region, fusion no longer provides net energy, because the
binding-energy curve has reached its maximum.

%%═══════════════════════════════════════════════════════════════════
\subsection{Fusion Power on Earth}
\label{subsec:lec11:fusion_power}
%%═══════════════════════════════════════════════════════════════════

Terrestrial fusion needs three things simultaneously:
\begin{enumerate}[nosep]
  \item sufficiently high temperature, so some particles enter the Gamow window;
  \item sufficiently high density, so collisions occur often enough;
  \item sufficiently long confinement time, so the plasma releases more fusion
        energy than the machine spends maintaining it.
\end{enumerate}

\dfn[dfn:lec11:fusion_q]{Fusion Gain}{%
The physical fusion gain is often written
\begin{equation}
  Q_{\rm phys}=\frac{\text{fusion energy released by the plasma}}
                    {\text{energy delivered to the plasma}}.
\end{equation}
Engineering gain compares useful electrical output with the total energy cost of
running the entire facility, including magnets, lasers, cryogenics, pumps,
control systems, and conversion losses.}

\paragraph{Magnetic confinement.}
In a tokamak or stellarator, a hot ionized plasma is confined by magnetic fields,
because no material wall can directly hold a $\sim10^8\,\mathrm K$ plasma. A
major design goal is to keep the plasma stable and confined long enough that
fusion heating becomes self-sustaining. ITER is the large international tokamak
project in France. Its goal is a large physical gain, not commercial electricity
production.

\paragraph{Inertial confinement.}
In laser inertial confinement, intense laser pulses compress a tiny capsule of
fuel for a very short time. The National Ignition Facility (NIF) is the main
example. Such experiments can reach physical break-even in a carefully defined
sense, but this is still far from an engineering power plant, because the driver
system and repetition-rate requirements are enormous.

\nt{Headlines about ``break-even'' are therefore ambiguous. A device may have
$Q_{\rm phys}>1$ for the plasma while still having engineering gain far below
unity for the full facility. A power plant needs the latter.}

%%═══════════════════════════════════════════════════════════════════
\subsection{Fission as a Collective Instability}
\label{subsec:lec11:fission_liquid_drop}
%%═══════════════════════════════════════════════════════════════════

\dfn[dfn:lec11:fission]{Fission}{%
Fission is the splitting of a heavy nucleus into two substantial fragments,
typically with several free neutrons and roughly $200\,\mathrm{MeV}$ of released
energy per fission event.}

Fission can be viewed as a very asymmetric relative of alpha decay: both require
separation of charged nuclear pieces through a Coulomb barrier. The difference is
that in alpha decay the emitted cluster is a compact, strongly bound
\nucleus{He}{4}; in fission there is no reason to imagine a preformed fragment
with $Z\sim30$--$60$ sitting inside the nucleus. Instead, the whole heavy nucleus
undergoes a collective deformation.

\subsubsection*{Liquid-Drop Energy Balance}
The relevant terms in the semi-empirical mass formula are the surface and
Coulomb terms,
\begin{equation}
  B(A,Z)=\cdots -a_s A^{2/3}-a_c\frac{Z^2}{A^{1/3}}+\cdots,
\end{equation}
where using $Z^2$ rather than $Z(Z-1)$ is sufficient for this qualitative
large-$Z$ argument. The surface term favours a sphere, because for fixed volume
a sphere minimizes surface area. The Coulomb term favours deformation, because
protons reduce their mutual repulsion by moving farther apart.

Model the nucleus as a slightly deformed spheroid with semiaxes
\begin{equation}
  a=R(1+\varepsilon),
  \qquad
  b=R(1+\varepsilon)^{-1/2},
\end{equation}
so that the volume is conserved:
\begin{equation}
  ab^2=R^3.
\end{equation}
For small $\varepsilon$, the surface area and Coulomb energy have expansions of
the form
\begin{align}
  S(\varepsilon)
  &=4\pi R^2\left(1+\frac{2}{5}\varepsilon^2+\cdots\right),
  \label{eq:lec11:surface_expansion}\\
  E_C(\varepsilon)
  &=E_C(0)\left(1-\frac{1}{5}\varepsilon^2+\cdots\right).
  \label{eq:lec11:coulomb_expansion}
\end{align}
There is no term linear in $\varepsilon$ because the sphere is an extremum.

The binding-energy change is therefore
\begin{equation}
  \Delta B
  =B(\varepsilon)-B(0)
  \simeq
  \left[-\frac{2}{5}a_sA^{2/3}
  +\frac{1}{5}a_cZ^2A^{-1/3}\right]\varepsilon^2.
  \label{eq:lec11:deltaB_deformation}
\end{equation}
The first term is negative: increasing surface area reduces binding. The second
term is positive: reducing Coulomb repulsion increases binding.

The spherical shape becomes unstable if a small deformation increases the
binding energy, i.e. if the coefficient of $\varepsilon^2$ is positive:
\begin{equation}
  a_cZ^2A^{-1/3}>2a_sA^{2/3}.
\end{equation}
Equivalently,
\begin{equation}
  \boxed{\frac{Z^2}{A}>\frac{2a_s}{a_c}\approx 47.}
  \label{eq:lec11:fissility_condition}
\end{equation}
This is the liquid-drop fissility criterion.

\paragraph{Physical intuition.}
For ordinary heavy nuclei, the surface term still wins near the spherical shape,
so there is a fission barrier. But as $Z^2/A$ grows, Coulomb repulsion becomes
more dangerous. Superheavy nuclei are therefore limited by fission instability,
although shell closures can produce enhanced stability in special regions.

%%═══════════════════════════════════════════════════════════════════
\subsection{Induced Fission and Fission Products}
\label{subsec:lec11:induced_fission}
%%═══════════════════════════════════════════════════════════════════

A neutron can induce fission by being captured and forming a highly excited
compound nucleus. A representative channel is
\begin{equation}
  \nucleus{U}{235}+n\longrightarrow \nucleus{U}{236}^{\ast}
  \longrightarrow \nucleus{Ba}{141}+\nucleus{Kr}{93}+2n+Q,
  \label{eq:lec11:u235_fission_example}
\end{equation}
with
\begin{equation}
  Q\sim 200\,\mathrm{MeV}.
\end{equation}
The exact fragments vary from event to event. The mass distribution is broad but
not centred at equal halves; for \nucleus{U}{235} thermal fission it has two
large peaks roughly near
\begin{equation}
  A\sim 95,
  \qquad
  A\sim 140.
\end{equation}
Thus fission usually gives one lighter and one heavier fragment, plus prompt
neutrons.

\paragraph{Prompt neutrons.}
The average number of prompt neutrons emitted per fission is denoted $\nu$ and
is typically
\begin{equation}
  \boxed{\nu\sim2\text{--}4.}
\end{equation}
For \nucleus{U}{235}, the average is about $2.4$--$2.5$. These neutrons are the
reason a chain reaction is possible: one neutron can trigger fission, but the
fission produces more than one neutron.

\subsubsection*{Energy Sharing Between Fragments}
In the centre-of-mass frame, the two large fragments recoil nearly back-to-back.
Neglecting the relatively small momentum carried by prompt neutrons, momentum
conservation gives
\begin{equation}
  p_1=p_2=p.
\end{equation}
Their kinetic energies are
\begin{equation}
  T_i=\frac{p^2}{2M_i}.
\end{equation}
Therefore
\begin{equation}
  \boxed{\frac{T_1}{T_2}=\frac{M_2}{M_1}.}
  \label{eq:lec11:fission_energy_sharing}
\end{equation}
The lighter fragment receives the larger kinetic energy. This kinetic energy is
quickly converted into heat in the surrounding material.

%%═══════════════════════════════════════════════════════════════════
\subsection{Why \texorpdfstring{$\nucleus{U}{235}$}{U-235} Fissions with Slow Neutrons but \texorpdfstring{$\nucleus{U}{238}$}{U-238} Does Not}
\label{subsec:lec11:u235_u238}
%%═══════════════════════════════════════════════════════════════════

A crucial point in reactor physics is that \nucleus{U}{235} has a large fission
cross section for thermal neutrons, whereas \nucleus{U}{238} does not. The
lecture explanation used pairing.

\paragraph{Pairing-energy argument.}
Uranium has $Z=92$, even. For \nucleus{U}{238}, the neutron number is
$N=146$, also even. Capturing a neutron gives $\nucleus{U}{239}^{\ast}$, with odd
$N=147$. This loses pairing energy relative to an even-neutron configuration.
For \nucleus{U}{235}, however, $N=143$ is odd. Capturing a neutron gives
$\nucleus{U}{236}^{\ast}$ with even $N=144$, gaining pairing energy. That extra
excitation helps the compound nucleus cross the fission barrier.

Because barrier penetration and fission probabilities depend very sensitively on
available excitation energy, a pairing-energy difference of order MeV causes
orders-of-magnitude differences in cross section.

\paragraph{Energy dependence.}
The thermal fission cross section of \nucleus{U}{235} is very large at slow
neutron energies and contains resonant structure at higher energies. For
\nucleus{U}{238}, fission becomes significant mainly for fast neutrons above a
threshold. Hence reactors that rely on thermal neutrons need a fissile isotope
such as \nucleus{U}{235} or \nucleus{Pu}{239}.

\nt{Slow neutrons are more effective partly because their de Broglie wavelength
is larger: $\lambda=h/p$. A larger wave character means a larger effective
interaction cross section. For many thermal-neutron reactions the cross section
roughly follows a $1/v$ trend away from resonances.}

%%═══════════════════════════════════════════════════════════════════
\subsection{Neutron Multiplication: The Four-Factor Formula}
\label{subsec:lec11:four_factor}
%%═══════════════════════════════════════════════════════════════════

A reactor is controlled by the average number of neutrons in one generation
relative to the previous generation. This is the multiplication factor.

\dfn[dfn:lec11:k]{Multiplication Factor}{%
The effective multiplication factor is
\begin{equation}
  k=\frac{\text{number of neutrons in generation }n+1}
          {\text{number of neutrons in generation }n}.
\end{equation}
If $k<1$ the chain reaction dies out; if $k=1$ it is critical and steady; if
$k>1$ it grows.}

For an ideal infinite reactor, no neutron leaks out. The multiplication factor is
then written
\begin{equation}
  \boxed{k_\infty=\eta\,\varepsilon\,p\,f.}
  \label{eq:lec11:four_factor}
\end{equation}
These four factors follow the life history of a neutron.

\begin{description}[leftmargin=2.8cm,style=nextline]
  \item[$\eta$ --- reproduction factor]
  Number of fast neutrons produced per thermal neutron absorbed in the fuel:
  \begin{equation}
    \eta=\nu\frac{\sigma_f}{\sigma_a}
    \simeq \nu\frac{\sigma_f}{\sigma_f+\sigma_\gamma}.
    \label{eq:lec11:eta}
  \end{equation}
  Here $\sigma_f$ is the fission cross section and $\sigma_a$ is the total
  absorption cross section in the fuel.

  \item[$\varepsilon$ --- fast fission factor]
  Correction for additional fissions caused by fast neutrons before they slow
  down. Usually $\varepsilon>1$.

  \item[$p$ --- resonance escape probability]
  Probability that a neutron slows from fast to thermal energies without being
  absorbed in resonance capture, especially by \nucleus{U}{238}.

  \item[$f$ --- thermal utilization factor]
  Probability that a thermal neutron is absorbed in the fuel rather than in the
  moderator, coolant, control material, or structural material.
\end{description}

\paragraph{Derivation as a neutron balance.}
Start with one thermal neutron absorbed in the fuel. It produces on average
$\eta$ fast neutrons. Fast fission increases this by $\varepsilon$. A fraction
$p$ survives slowing down without resonance absorption. Finally, a fraction $f$
of the thermalized neutrons is absorbed in the fuel where it can continue the
chain. Multiplying these independent survival/production factors gives
Eq.~\eqref{eq:lec11:four_factor}.

%%═══════════════════════════════════════════════════════════════════
\subsection{Finite Reactors: Leakage and the Six-Factor Formula}
\label{subsec:lec11:six_factor}
%%═══════════════════════════════════════════════════════════════════

Real reactors are finite, so neutrons can escape. The infinite multiplication
factor must be multiplied by non-leakage probabilities. In the lecture notation,
if $N_f$ and $N_t$ are leakage probabilities for fast and thermal neutrons, then
\begin{equation}
  \boxed{k_{\rm eff}
  =\eta\,\varepsilon\,p\,f\,(1-N_f)(1-N_t).}
  \label{eq:lec11:six_factor}
\end{equation}
Equivalently, many textbooks write
\begin{equation}
  k_{\rm eff}=\eta\,\varepsilon\,p\,f\,P_f\,P_t,
\end{equation}
where $P_f$ and $P_t$ are fast and thermal non-leakage probabilities.

\paragraph{Critical size intuition.}
Large systems leak a smaller fraction of their neutrons because volume grows as
$R^3$ while surface area grows as $R^2$. This is why a sufficiently small piece
of fissile material can be subcritical even if the same composition in a larger
geometry would be critical.

%%═══════════════════════════════════════════════════════════════════
\subsection{Criticality, Prompt Neutrons, and Delayed Neutrons}
\label{subsec:lec11:criticality_delayed}
%%═══════════════════════════════════════════════════════════════════

If each generation takes a characteristic time $\tau_g$, then after $n$
generations the neutron population is approximately
\begin{equation}
  N_n=N_0k^n,
  \qquad
  n=\frac{t}{\tau_g}.
\end{equation}
Thus
\begin{equation}
  \boxed{N(t)=N_0\,k^{t/\tau_g}
  =N_0\exp\left(\frac{\ln k}{\tau_g}t\right).}
  \label{eq:lec11:neutron_growth}
\end{equation}
If $\tau_g\sim10^{-3}\,\mathrm s$ and $k=1.1$, then after one second
\begin{equation}
  N(1\,\mathrm s)=N_0(1.1)^{1000},
\end{equation}
which is catastrophically large. A reactor cannot be safely controlled if it is
supercritical on prompt neutrons alone.

The solution is delayed neutrons. Some neutron-rich fission fragments undergo
$\beta^-$ decay to an excited daughter state, and the daughter excitation lies
above the neutron separation energy $S_n$. The daughter can then emit a neutron:
\begin{equation}
  (A,Z)\xrightarrow{\beta^-}(A,Z+1)^\ast
  \longrightarrow (A-1,Z+1)+n.
\end{equation}
These \emph{beta-delayed neutrons} appear on time scales of seconds rather than
milliseconds.

\thm[th:lec11:delayed_neutrons]{Why Delayed Neutrons Make Reactors Controllable}{%
A power reactor is designed so that it is subcritical on prompt neutrons alone,
but critical when delayed neutrons are included. Then the relevant control time
scale is seconds, not milliseconds, allowing mechanical control rods and thermal
feedback to regulate the chain reaction.}

%%═══════════════════════════════════════════════════════════════════
\subsection{Main Reactor Ingredients}
\label{subsec:lec11:reactor_ingredients}
%%═══════════════════════════════════════════════════════════════════

A thermal reactor needs four basic ingredients.

\begin{enumerate}[label=\textbf{\arabic*.},leftmargin=1.2cm]
  \item \textbf{Fissile material.} A fuel such as \nucleus{U}{235} or
        \nucleus{Pu}{239} that can fission efficiently after thermal-neutron
        absorption.
  \item \textbf{Moderator.} A material that slows fast neutrons to thermal
        energies. Efficient slowing requires elastic collisions with nuclei of
        comparable mass, so hydrogen is excellent.
  \item \textbf{Control rods.} Neutron absorbers such as cadmium, boron, or
        hafnium. Inserting rods lowers $k$; withdrawing rods raises $k$.
  \item \textbf{Coolant.} A circulating material that removes heat from the
        core and transfers it to a heat exchanger or turbine system.
\end{enumerate}

\subsubsection*{Why Hydrogen Moderates Well}
For an elastic collision between a neutron of mass $m$ and a stationary target
of mass $M$, the maximum fractional energy transfer is
\begin{equation}
  \boxed{\frac{\Delta E_{\max}}{E}
  =\frac{4mM}{(m+M)^2}.}
  \label{eq:lec11:elastic_transfer}
\end{equation}
This is maximal when $M=m$. Therefore protons in ordinary water are extremely
good neutron moderators. Concrete also shields neutrons well partly because it
contains water/hydrogen.

\paragraph{Light water.}
Ordinary water moderates well, is cheap, and is an excellent heat-transfer fluid.
Its disadvantage is that hydrogen also absorbs some neutrons, reducing $f$.
Light-water reactors therefore usually require enriched uranium.

\paragraph{Heavy water.}
Heavy water, $\mathrm{D_2O}$, moderates somewhat less efficiently per collision
because deuterium is heavier than a neutron, but it absorbs far fewer neutrons.
This allows CANDU-type reactors to run on natural uranium, at the price of using
expensive heavy water.

\paragraph{Graphite.}
Graphite is less effective per collision than hydrogen but has low absorption and
is mechanically convenient. Graphite-moderated reactors need a separate coolant,
such as water or gas.

%%═══════════════════════════════════════════════════════════════════
\subsection{Heat Engines and Reactor Types}
\label{subsec:lec11:reactor_types}
%%═══════════════════════════════════════════════════════════════════

Nuclear reactors do not directly turn nuclear energy into electricity. They make
heat. The heat is used to run a thermodynamic cycle, usually by producing steam
that drives turbines.

The ideal upper bound on heat-engine efficiency is the Carnot efficiency,
\begin{equation}
  \boxed{\eta_{\rm Carnot}=1-\frac{T_c}{T_h}.}
  \label{eq:lec11:carnot}
\end{equation}
Thus a hotter reactor can in principle convert heat to electricity more
efficiently. Ordinary water at atmospheric pressure boils at
$100^\circ\mathrm C$, too low for high efficiency, so many reactors operate water
under high pressure.

\paragraph{Pressurized-water reactor (PWR).}
The core water is kept at high pressure so it does not boil, with typical
operating temperatures of order $300^\circ\mathrm C$. It transfers heat to a
secondary water loop that makes steam for the turbine. This keeps the most
radioactive water inside the primary loop.

\paragraph{Boiling-water reactor (BWR).}
Water boils in the core and the steam goes more directly to the turbine system.
The design is simpler in some ways, but the steam loop is more radioactive.

%%═══════════════════════════════════════════════════════════════════
\subsection{Enrichment, Natural Uranium, and Oklo}
\label{subsec:lec11:enrichment_oklo}
%%═══════════════════════════════════════════════════════════════════

Natural uranium today is mostly \nucleus{U}{238}, with only about
$0.7\%$ \nucleus{U}{235}. Light-water reactors usually require enrichment to a
few percent \nucleus{U}{235}. A common enrichment route uses uranium hexafluoride,
$\mathrm{UF_6}$, in gas centrifuges. Because molecules containing
\nucleus{U}{235} are slightly lighter than those containing \nucleus{U}{238}, a
rapidly rotating centrifuge produces a tiny isotope separation. Cascading many
centrifuges gives useful enrichment.

Historically, gaseous diffusion was also used: the diffusion rate depends weakly
on molecular mass, so repeated diffusion stages can enrich \nucleus{U}{235}, but
very inefficiently.

\paragraph{Oklo: a natural reactor.}
A famous natural fission reactor operated in Oklo, Gabon, roughly two billion
years ago. The key fact is that \nucleus{U}{235} has a shorter half-life than
\nucleus{U}{238}, so the natural \nucleus{U}{235} fraction was higher in the
past. In special geological conditions, with water acting as a moderator, a
natural uranium deposit could become critical. Later measurements found uranium
depleted in \nucleus{U}{235}, revealing that fission had already occurred there.

%%═══════════════════════════════════════════════════════════════════
\subsection{Waste, Decay Heat, and the Nuclear-Power Debate}
\label{subsec:lec11:waste}
%%═══════════════════════════════════════════════════════════════════

Fission products and neutron-activation products are radioactive. Even after the
chain reaction stops, the fuel continues to release decay heat. This is why spent
fuel is first stored in water pools: the water removes heat and provides
radiation shielding. After shorter-lived isotopes decay, fuel can be moved to dry
casks or other storage systems.

\paragraph{Long time scales.}
The difficult waste problem is not the first few seconds or days; it is the
combination of radiotoxicity, heat production, chemical toxicity, and very long
half-lives. A natural benchmark is: how long until the waste is no more
radioactive than the uranium ore originally mined? The answer can be of order
$10^4$--$10^5$ years, depending on the fuel cycle and what isotopes are included.
This motivates geological repositories in stable rock formations.

\paragraph{Why the debate is hard.}
Nuclear fission produces large amounts of low-carbon electricity, but it raises
hard questions: accident risk, cost of construction and decommissioning,
long-term waste stewardship, proliferation concerns, and public trust. Newer
reactor concepts attempt to improve safety, burn more of the fuel, use waste as
fuel, or reduce long-lived waste. Whether those designs become economical and
widely deployable is still an open technological and social question.

%%═══════════════════════════════════════════════════════════════════
\subsection*{Conceptual Summary}
%%═══════════════════════════════════════════════════════════════════

\begin{itemize}
  \item Fusion and fission both release energy by moving nuclear matter toward
        larger total binding energy.
  \item Fusion is limited by Coulomb-barrier penetration; stars exploit the
        Gamow window plus gravity and enormous time.
  \item Terrestrial fusion must distinguish physical plasma gain from full
        engineering gain.
  \item Fission of heavy nuclei is a competition between surface energy, which
        favours a sphere, and Coulomb energy, which favours elongation and
        splitting.
  \item \nucleus{U}{235} is fissile with thermal neutrons because neutron capture
        gives favourable excitation/pairing in $\nucleus{U}{236}^{\ast}$;
        \nucleus{U}{238} mainly fissions with fast neutrons.
  \item A reactor is a neutron economy. The four-factor formula tracks
        production, fast fission, resonance survival, and thermal utilization;
        the six-factor formula adds leakage.
  \item Safe power reactors operate with $k_{\rm eff}=1$ and rely on delayed
        neutrons plus control rods and feedback to make the dynamics slow enough
        to control.
\end{itemize}
